\clearpage
\begin{summarybox}
\textbf{C092\quad 共焦椭圆与双曲线在公共点处正交：快记／快懂／快用}

\textbf{快记（统一记号＋两条结论）：}椭圆用小写 $a,b$，双曲线用大写 $A,B$。
当 $a>c>A>0$，椭圆与双曲线有相同的两个焦点时，在所有公共点处
\[
\boxed{t_E\perp t_H},\qquad \boxed{a^2-A^2=b^2+B^2}.
\]
\textbf{快懂（抓“和、差”对应的方向）：}
\[
PF_1+PF_2=2a,\quad|PF_1-PF_2|=2A,
\]
\[
\mathbf n_E\parallel\mathbf u_1+\mathbf u_2,\quad
\mathbf n_H\parallel\mathbf u_1-\mathbf u_2,
\quad |\mathbf u_1|=|\mathbf u_2|=1.
\]
\[
\boxed{(\mathbf u_1+\mathbf u_2)\cdot
(\mathbf u_1-\mathbf u_2)=1-1=0}.
\]
\textbf{一句几何解释：}椭圆的切线沿焦角\emph{外}角平分线，
双曲线的切线沿焦角\emph{内}角平分线；内、外角平分线垂直。

\textbf{为何有平方关系：}
$c_E^2=a^2-b^2$，$c_H^2=A^2+B^2$；共焦使两式相等。
同轴标准式下也可逆用。\textbf{公共点本身不足以推出平方关系。}
\[\boxed{x_0^2=\frac{a^2A^2}{c^2}},\qquad
\boxed{y_0^2=\frac{b^2B^2}{c^2}}.\]

\textbf{快用（考场3步）：}
\begin{enumerate}
\item 看见“椭圆＋双曲线”，先查 $a^2-b^2=A^2+B^2$，不能仅看共轴；
\item 若共焦且在同一点相交，\emph{直接调用正交定理}；
\item 若要求证明，优先单位向量和差点积；需具体切线则用隐函数求导。
\end{enumerate}
\textbf{最小反例：}非共焦的
$E:x^2/16+y^2/12=1$、$H_1:x^2/2-y^2/9=1$
同过 $(2,3)$，但切线斜率积 $-3/2\ne-1$。

\textbf{关联：}C090 焦角与离心率关系；C092 切线交角恒为 $90^\circ$。
\textbf{迁移线索：}两段距离的和、差保持不变的曲线、
角平分线、等长向量的和与差。
\end{summarybox}
